\chapter{Vận hành stack: Cẩm nang gỡ lỗi thực địa}
\label{ch:operations}

Đây là chương để đọc lại khi có một cluster hỏng trước mặt bạn. Nó
tập hợp quy trình gỡ lỗi mà các chương trước đã tích lũy, sắp xếp
theo cách sự cố thật diễn ra.

\section{Bốn câu hỏi}

Gần như mọi sự cố Kubernetes đều được định vị bằng cách hỏi, theo thứ
tự:

\begin{enumerate}
  \item \textbf{Pod có đang chạy không?} \code{kubectl -n ts get pods}
  --- đọc cột STATUS và RESTARTS; bảng phân loại bên dưới.
  \item \textbf{Nền tảng nói gì về nó?}
  \code{kubectl -n ts describe pod <pod>} --- mục Events ở cuối là
  kubelet và scheduler kể lại các quyết định của chúng: kéo image,
  probe thất bại, OOM kill, từ chối lập lịch.
  \item \textbf{Ứng dụng nói gì?} \code{kubectl -n ts logs
  <pod>} (thêm \code{--previous} cho container vừa chết --- cờ hay bị
  quên nhất trong ứng phó sự cố).
  \item \textbf{Tôi có tái hiện được góc nhìn của nó về thế giới
  không?} \code{kubectl -n ts exec -it <pod> -- sh} rồi thử kết nối
  mà ứng dụng khẳng định là đang thất bại.
\end{enumerate}

\section{Bảng phân loại trạng thái}

\begin{table}[h]
\centering\small
\begin{tabular}{@{}lp{6.6cm}@{}}
\toprule
Trạng thái & Ý nghĩa và bước đi đầu tiên \\
\midrule
\code{Pending} & không lập lịch được: requests vượt quá tài nguyên
còn trống, hoặc PVC chưa được gắn. \code{describe} cho biết là gì. \\
\code{ImagePullBackOff} & registry không truy cập được, gõ sai tên,
hoặc bẫy HTTPS-so-với-HTTP của Chương~13. \\
\code{CrashLoopBackOff} & tiến trình thoát ngay sau khi khởi động;
thời gian chờ nhân đôi giữa các lần thử. \code{logs --previous}. \\
\code{OOMKilled} (exit 137) & chạm giới hạn bộ nhớ; không có trace
Java nào --- cú kill từ kernel của Chương~17. Tăng limit hoặc thu nhỏ
heap. \\
\code{Running 0/1} & còn sống nhưng readiness thất bại: bình thường
khi khởi động; là vấn đề khi kéo dài --- thử chính URL của probe. \\
\code{Completed}/\code{Error} & cho Job; với pod service,
\code{Error} nghĩa là tiến trình đã thoát với mã khác 0 một lần. \\
\bottomrule
\end{tabular}
\caption{Cột STATUS đang cố nói với bạn điều gì.}
\end{table}

\begin{handson}[cố ý làm hỏng --- ba bài diễn tập]
Phá hoại có chủ đích, vào một ngày yên tĩnh, dạy nhiều hơn bất kỳ sự
cố nào:
\begin{lstlisting}[language=bash]
# 1. wrong image tag -> watch ImagePullBackOff, then roll back
kubectl -n ts set image deployment/course-service \
    course-service=localhost:5000/course-service:nope
kubectl -n ts rollout undo deployment/course-service

# 2. scale the database down -> watch dependents degrade honestly
kubectl -n ts scale statefulset/postgres-course --replicas=0
kubectl -n ts get pods -w     # course-service readiness reacts
kubectl -n ts scale statefulset/postgres-course --replicas=1

# 3. kill Eureka -> verify existing routes survive on client caches
kubectl -n ts delete pod -l app=discovery-server
\end{lstlisting}
Mỗi bài diễn tập ghép lý thuyết của một chương với hành vi quan sát
được của nó: 13 (đặt tên), 17--18 (probe và endpoint), 3 (cache phía
client).
\end{handson}

\section{Đọc thế giới của một service từ bên trong}

Khi log khẳng định kết nối thất bại, hãy kiểm chứng lời khẳng định đó
từ chính vị trí mạng của pod --- DNS và khả năng truy cập khác nhau
giữa máy bạn và pod:

\begin{lstlisting}[language=bash]
kubectl -n ts exec -it deploy/course-service -- sh
wget -qO- http://config-server:8888/actuator/health   # DNS + service up?
nc -zv postgres-course 5432                           # TCP reachable?
env | grep SPRING                                     # what was injected?
\end{lstlisting}

Lệnh cuối cùng đó dàn xếp tranh cãi rất nhanh: nó in ra cấu hình mà
pod \emph{thật sự} nhận được, hơn hẳn việc đọc lại manifest.

\section{Thói quen giữ cho cuối tuần được rảnh}

\begin{description}
  \item[Sau mỗi lần deploy] pipeline đã kiểm chứng rollout rồi; dù
  vậy vẫn liếc qua \code{kubectl -n ts get pods} --- RESTARTS nhích
  dần lên là sự cố ngày mai đang tự báo trước.
  \item[Hằng tuần] \code{kubectl top pods -n ts} (cần metrics-server,
  đi kèm trong k3s): so sánh mức dùng thực tế với requests của
  Chương~17 và sửa lại phần hư cấu. Kiểm tra đĩa: \code{df -h} trên
  máy chủ --- registry và file log lớn dần; \code{registry:2} thỉnh
  thoảng cần dọn rác.
  \item[Trước buổi demo] \emph{đừng build} --- hộp phần cứng của
  Chương~5: hai nhân nghĩa là một lần build Maven và buổi demo tranh
  nhau đúng phần CPU đang giữ cho readiness probe trả lời.
\end{description}

\begin{hardware}
ufw trên máy chủ phải miễn trừ các CIDR của pod và service
(\code{10.42.0.0/16}, \code{10.43.0.0/16}) --- chặn tường lửa chúng
sẽ phá DNS và lưu lượng pod-tới-pod theo những cách giả dạng lỗi ứng
dụng: timeout chập chờn, tên lúc phân giải được lúc không. Khi mạng
trên k3s xảy ra chuyện \emph{kỳ lạ}, hãy nghi ngờ tường lửa máy chủ
trước ứng dụng. Và k3s chỉ đọc \code{registries.yaml} lúc khởi động
--- sau khi sửa, \code{systemctl restart k3s}, lệnh này khởi động lại
mọi pod trong chốc lát: đừng làm giữa buổi demo.
\end{hardware}

\begin{pitfall}[gỡ lỗi sai tầng]
Thói quen đắt giá nhất khi gỡ lỗi Kubernetes là bắt đầu từ nơi triệu
chứng xuất hiện thay vì nơi nguyên nhân sống. Một lỗi 502 trên trình
duyệt có thể là Traefik (không có endpoint), gateway (Eureka lạc
hậu), service (readiness thất bại), cơ sở dữ liệu (pod bị đuổi), hay
đĩa (đầy). Bốn câu hỏi tồn tại để đi qua các tầng theo thứ tự; hãy
cưỡng lại cám dỗ mở log ứng dụng trước chỉ vì đó là tầng quen thuộc
nhất. \code{describe} trước \code{logs} --- nền tảng thường đã biết
rồi.
\end{pitfall}

\begin{quiz}
\begin{enumerate}
  \item Sắp xếp bốn câu hỏi và nêu lệnh cho mỗi câu. Câu nào dùng một
  cờ mà mọi người hay quên, và khi nào cờ đó quan trọng?
  \item \code{CrashLoopBackOff} so với \code{Running 0/1}: các tầng
  khác nhau đang thất bại --- khái niệm probe nào phân biệt chúng?
  \item Một service khẳng định không tới được Postgres; từ laptop
  của bạn Postgres trả lời. Bước chẩn đoán tiếp theo là gì và vì sao
  thành công trên laptop của bạn không liên quan?
  \item Vì sao một thay đổi tường lửa máy chủ có thể giả dạng lỗi ứng
  dụng trên k3s, và hai CIDR nào phải luôn được miễn trừ?
\end{enumerate}
\end{quiz}
